\section{Discussion}

The institutional-friction results show that the ranking between hybrid and top-down-only governance depends on the scenario and on the practical limits of oversight. Community-irrigation settings favour hybrid governance most consistently. Constrained forest co-management and the higher-pressure regulated-fishery setting favour top-down-only governance.

These results narrow the scope of the earlier Harvest finding. The confirmatory architecture matrix still shows a strong result for hybrid governance under search-based pressure. After scenario structure and oversight frictions are introduced, that ordering no longer holds uniformly. The extension therefore changes the interpretation of the architecture comparison from a single ranking to a family of scenario-specific rankings.

The scenarios remain archetypal and are not calibrated to particular cases. The results support comparison across governance settings, but they should not be read as predictions for any one fishery, irrigation system, or forest governance arrangement. Each scenario family is represented by a single benchmark parameterization, and the institutional-friction extension fixes the attacker family to the strongest non-LLM search-based pressure instead of re-running the full attacker ladder within every new scenario.

Within those limits, the cumulative evidence is substantial. Fishery Commons ranks central interventions under repeated strategic turnover. Harvest Commons distinguishes local-only, central-only, and hybrid governance in a richer common-pool setting. The confirmatory architecture matrix, the capability ladder, the welfare-incidence analysis, and the institutional-friction extension together support a comparative claim about governance architecture under strategic pressure.
